\chapter{2009年北京卷（文）}

\section{单选题}

\begin{problem}
设集合 \(A = \left\{x \mid -\frac{1}{2} < x < 2\right\}\)，\(B = \{x \mid x^2 \leqslant 1\}\)，则 \(A \cup B =\)（\quad）

\choices
    {\(\{x\mid -1\leqslant x < 2\}\)}
    {\(\left\{x\mid -\frac{1}{2} < x\leqslant 1\right\}\)}
    {\( \{x \mid x < 2\} \)}
    {\(\{x\mid 1\leqslant x < 2\}\)}
\begin{answer}
A
\end{answer}

\begin{solution}
由 \(x^2\leqslant 1\) 得 \(-1\leqslant x\leqslant 1\)，所以 \(B=\left[-1,1\right]\). 又 \(A=\left(-\frac{1}{2},2\right)\)，两集合的并集为 \(A\cup B=\left[-1,2\right)\)，故选 A.
\end{solution}
\end{problem}

\begin{problem}
已知向量 \(\bs{a} = (1,0)\)，\(\bs{b} = (0,1)\)，\(\bs{c} = k\bs{a} + \bs{b}\)（\(k \in \R\)），\(\bs{d} = \bs{a} - \bs{b}\)，如果 \(\bs{c} \parallel \bs{d}\)，那么（\quad）

\choices
    {\(k = 1\) 且 \(\bs{c}\) 与 \(\bs{d}\) 同向}
    {\(k = 1\) 且 \(\bs{c}\) 与 \(\bs{d}\) 反向}
    {\(k = -1\) 且 \(\bs{c}\) 与 \(\bs{d}\) 同向}
    {\(k = -1\) 且 \(\bs{c}\) 与 \(\bs{d}\) 反向}
\begin{answer}
D
\end{answer}

\begin{solution}
由题设，\(\bs{c}=(k,1)\)，\(\bs{d}=(1,-1)\). 两向量平行的充要条件是行列式为零，即 \(k\cdot(-1)-1\cdot1=0\)，解得 \(k=-1\). 此时 \(\bs{c}=(-1,1)=-\bs{d}\)，所以两向量反向，故选 D.
\end{solution}
\end{problem}

\begin{problem}
若 \((1 + \sqrt{2})^4 = a + b\sqrt{2}\)（\(a\)，\(b\) 为有理数），则 \(a + b =\)（\quad）

\choices
    {\(33\)}
    {\(29\)}
    {\(23\)}
    {\(19\)}
\begin{answer}
B
\end{answer}

\begin{solution}
先平方得 \(\left(1+\sqrt{2}\right)^2=3+2\sqrt{2}\)，因此 \(\left(1+\sqrt{2}\right)^4=\left(3+2\sqrt{2}\right)^2=17+12\sqrt{2}\). 对照 \(a+b\sqrt{2}\)，得 \(a=17\)，\(b=12\)，所以 \(a+b=29\)，故选 B.
\end{solution}
\end{problem}

\begin{problem}
为了得到函数 \( y = \lg \frac{x + 3}{10} \) 的图象，只需把函数 \( y = \lg x \) 的图象上所有的点（\quad）

\choices
    {向左平移 \(3\) 个单位长度，再向上平移 \(1\) 个单位长度}
    {向右平移 \(3\) 个单位长度，再向上平移 \(1\) 个单位长度}
    {向左平移 \(3\) 个单位长度，再向下平移 \(1\) 个单位长度}
    {向右平移 \(3\) 个单位长度，再向下平移 \(1\) 个单位长度}
\begin{answer}
C
\end{answer}

\begin{solution}
将 \(y=\lg x\) 的图象向左平移 \(3\) 个单位长度，得到 \(y=\lg(x+3)\)；再向下平移 \(1\) 个单位长度，得到 \(y=\lg(x+3)-1=\lg\frac{x+3}{10}\). 因此应选 C.
\end{solution}
\end{problem}

\begin{problem}
用数字 1， 2， 3， 4， 5 组成的无重复数字的四位偶数的个数为（\quad）

\choices
    {\(8\)}
    {\(24\)}
    {\(48\)}
    {\(120\)}
\begin{answer}
C
\end{answer}

\begin{solution}
四位偶数的个位只能取 \(2\) 或 \(4\)，有 \(2\) 种取法. 个位确定后，其余三位从剩下的 \(4\) 个数字中选取并排列，方法数为 \(\mathrm A_4^3=4\cdot3\cdot2=24\). 所以共有 \(2\times24=48\) 个，故选 C.
\end{solution}
\end{problem}

\begin{problem}
“ \( a=\frac{\pi}{6} \) ”是“ \( \cos2a=\frac{1}{2} \) ”的（\quad）

\choices
    {充分而不必要条件}
    {必要而不充分条件}
    {充分必要条件}
    {既不充分也不必要条件}
\begin{answer}
A
\end{answer}

\begin{solution}
当 \(a=\frac{\pi}{6}\) 时，\(\cos2a=\cos\frac{\pi}{3}=\frac{1}{2}\)，所以前一个条件能推出后一个条件. 反之，\(\cos2a=\frac{1}{2}\) 的解为 \(2a=2k\pi\pm\frac{\pi}{3}\)，即 \(a=k\pi\pm\frac{\pi}{6}\)（\(k\in\Z\)），不一定有 \(a=\frac{\pi}{6}\). 因此前者是后者的充分而不必要条件，故选 A.
\end{solution}
\end{problem}

\begin{problem}
若正四棱柱 \(ABCD - A_{1}B_{1}C_{1}D_{1}\) 的底面边长为 \(1\)，\(AB_{1}\) 与底面 \(ABCD\) 成 \(60^{\circ}\) 角，则 \(A_{1}C_{1}\) 到底面 \(ABCD\) 的距离为（\quad）

\choices
    {\(\frac{\sqrt{3}}{3}\)}
    {\(1\)}
    {\(\sqrt{2}\)}
    {\(\sqrt{3}\)}
\begin{answer}
D
\end{answer}

\begin{solution}
设正四棱柱的高为 \(h=BB_1\). 直线 \(AB_1\) 在底面上的射影为 \(AB\)，所以 \(\tan60^{\circ}=\frac{BB_1}{AB}=h\)，从而 \(h=\sqrt{3}\). 由于 \(A_1C_1\) 在与底面平行的上底面内，\(A_1C_1\) 到底面 \(ABCD\) 的距离就是棱柱的高 \(\sqrt{3}\)，故选 D.
\end{solution}
\end{problem}

\begin{problem}
设 \(D\) 是正 \(\triangle P_{1}P_{2}P_{3}\) 及其内部的点构成的集合，点 \(P_{0}\) 是 \(\triangle P_{1}P_{2}P_{3}\) 的中心. 若集合 \(S = \{P \mid P \in D, |PP_{0}| \leqslant |PP_{i}|, i = 1, 2, 3\}\)，则集合 \(S\) 表示的平面区域是（\quad）

\choices
    {三角形区域}
    {四边形区域}
    {五边形区域}
    {六边形区域}
\begin{answer}
D
\end{answer}

\begin{solution}
设正三角形的中心为 \(O=P_0\)，其外接圆半径为 \(R\). 对任意顶点 \(P_i\)，条件 \(\lvert PO\rvert\leqslant\lvert PP_i\rvert\) 等价于 \(\lvert\overrightarrow{OP}\rvert^2\leqslant\lvert\overrightarrow{OP}-\overrightarrow{OP_i}\rvert^2\)，即 \(\overrightarrow{OP}\mathbin{\cdot}\overrightarrow{OP_i}\leqslant\frac{R^2}{2}\). 因而每个条件确定一个以 \(OP_i\) 的垂直平分线为边界、且含有中心 \(O\) 的半平面.

三条边界线分别截去正三角形的三个顶角，每条边界线在三角形内截出一条线段；剩余区域的边界由原三角形的三段边和这三条截线组成，共有六条边，所以 \(S\) 表示六边形区域，故选 D.
\end{solution}
\end{problem}

\section{填空题}

\begin{problem}
若 \(\sin \theta = -\frac{4}{5}\)，\(\tan \theta > 0\)，则 \(\cos \theta =\fillinblank{}\).
\begin{answer}
\(-\frac{3}{5}\).
\end{answer}

\begin{solution}
因为 \(\tan\theta=\frac{\sin\theta}{\cos\theta}>0\)，且 \(\sin\theta<0\)，所以 \(\cos\theta<0\). 又由同角三角函数的基本关系，\(\cos^2\theta=1-\sin^2\theta=1-\frac{16}{25}=\frac{9}{25}\). 因此 \(\cos\theta=-\frac{3}{5}\).
\end{solution}
\end{problem}

\begin{problem}
若数列 \(\{a_{n}\}\) 满足： \(a_{1} = 1\)，\(a_{n+1} = 2a_{n}\)（\(n \in \N^*\)），则 \(a_{5} =\fillinblank{}\)；前 8 项的和 \(S_{8} =\fillinblank{}\). （用数字作答）
\begin{answer}
\(16\)，\(255\).
\end{answer}

\begin{solution}
由递推关系可得 \(a_n=2^{n-1}a_1=2^{n-1}\)，所以 \(a_5=2^4=16\). 前 8 项构成首项为 \(1\)、公比为 \(2\) 的等比数列，故 \(S_8=\frac{2^8-1}{2-1}=255\).
\end{solution}
\end{problem}

\begin{problem}
若实数 \(x\)，\(y\) 满足 \(\begin{cases}
x + y - 2 \geqslant 0,\\
x \leqslant 4,\\
y \leqslant 5.
\end{cases}\) 则 \(s = x + y\) 的最大值为\fillinblank{}.
\begin{answer}
\(9\).
\end{answer}

\begin{solution}
由 \(x\leqslant4\) 和 \(y\leqslant5\)，得 \(s=x+y\leqslant4+5=9\). 取 \(x=4\)，\(y=5\) 时，\(x+y-2=7\geqslant0\)，满足全部条件，且 \(s=9\). 所以 \(s\) 的最大值为 \(9\).
\end{solution}
\end{problem}

\begin{problem}
已知函数 \( f(x) = \begin{cases}
3^x, & x \leqslant 1,\\
-x, & x > 1.
\end{cases} \) 若 \( f(x) = 2 \)，则 \( x = \)\fillinblank{}.
\begin{answer}
\(\log_{3}2\).
\end{answer}

\begin{solution}
当 \(x\leqslant1\) 时，\(3^x=2\)，解得 \(x=\log_3 2\). 因为 \(1<2<3\)，所以 \(0<\log_3 2<1\)，满足该分支的定义域. 当 \(x>1\) 时，\(-x=2\) 得 \(x=-2\)，不满足 \(x>1\). 因此唯一解为 \(x=\log_3 2\).
\end{solution}
\end{problem}

\begin{problem}
椭圆 \(\frac{x^2}{9} + \frac{y^2}{2} = 1\) 的焦点为 \(F_1\)，\(F_2\)，点 \(P\) 在椭圆上，若 \(|PF_1| = 4\)，则 \(|PF_2| =\fillinblank{}\)；\(\angle F_1PF_2\) 的大小为\fillinblank{}.
\begin{answer}
\(2\)，\(120^{\circ}\).
\end{answer}

\begin{solution}
椭圆的长半轴为 \(a=3\)，短半轴为 \(b=\sqrt{2}\)，所以半焦距 \(c=\sqrt{a^2-b^2}=\sqrt7\)，从而 \(|F_1F_2|=2\sqrt7\). 由椭圆定义，\(|PF_1|+|PF_2|=2a=6\)，故 \(|PF_2|=6-4=2\). 在三角形 \(F_1PF_2\) 中，由余弦定理，\(\cos\angle F_1PF_2=\frac{4^2+2^2-(2\sqrt7)^2}{2\cdot4\cdot2}=-\frac12\). 因此 \(\angle F_1PF_2=120^{\circ}\).
\end{solution}
\end{problem}

\begin{problem}
设 \(A\) 是整数集的一个非空子集，对于 \(k \in A\)，如果 \(k - 1 \notin A\) 且 \(k + 1 \notin A\)，那么 \(k\) 是 \(A\) 的一个“孤立元”，给定 \(S = \{1, 2, 3, 4, 5, 6, 7, 8\}\)，由 \(S\) 的 3 个元素构成的所有集合中，不含“孤立元”的集合共有\fillinblank{} 个.
\begin{answer}
\(6\).
\end{answer}

\begin{solution}
把一个由 \(S\) 的元素构成的集合按连续整数分成若干段. 不含孤立元意味着每一段的长度都至少为 \(2\). 该集合共有 \(3\) 个元素，因此只能是一段长度为 \(3\) 的连续整数，即 \(\{1,2,3\}\)、\(\{2,3,4\}\)、\(\{3,4,5\}\)、\(\{4,5,6\}\)、\(\{5,6,7\}\)、\(\{6,7,8\}\). 共 \(6\) 个.
\end{solution}
\end{problem}

\section{解答题}

\begin{problem}
已知函数 \( f(x) = 2\sin (\pi - x)\cos x \).
\begin{enumerate}
\item 求 \( f(x) \) 的最小正周期；
\item 求 \( f(x) \) 在区间 \( \left[-\frac{\pi}{6}, \frac{\pi}{2}\right] \) 上的最大值和最小值.
\end{enumerate}

\begin{solution}
\begin{enumerate}
\item 因为 \(\sin(\pi-x)=\sin x\)，所以
\(f(x)=2\sin x\cos x=\sin 2x\).
函数 \(\sin 2x\) 的最小正周期 \(T\) 满足 \(2T=2\pi\)，故 \(T=\pi\).

\item 当 \(x\in\left[-\frac{\pi}{6},\frac{\pi}{2}\right]\) 时，
\(-\frac{\pi}{3}\leqslant 2x\leqslant\pi\).
在区间 \(\left[-\frac{\pi}{3},\pi\right]\) 上，\(\sin 2x\) 在 \(2x=\frac{\pi}{2}\) 时取到最大值 \(1\)，且在左端点取到最小值 \(\sin\left(-\frac{\pi}{3}\right)=-\frac{\sqrt{3}}{2}\).
因此，当 \(x=\frac{\pi}{4}\) 时，\(f(x)\) 取得最大值 \(1\)；当 \(x=-\frac{\pi}{6}\) 时，\(f(x)\) 取得最小值 \(-\frac{\sqrt{3}}{2}\).
\end{enumerate}
\end{solution}
\end{problem}

\begin{problem}
如图，四棱锥 \(P - ABCD\) 的底面是正方形，\(PD \perp ABCD\)，点 \(E\) 在棱 \(PB\) 上.
\begin{enumerate}
\item 求证：平面 \( AEC \perp \) 平面 \(PDB\)；
\item 当 \(PD = \sqrt{2} AB\) 且 \(E\) 为 \(PB\) 的中点时，求 \(AE\) 与平面 \(PDB\) 所成的角的大小.
\end{enumerate}

\texfigure[max width=0.38\linewidth]{img_repaint/2009/beijing_liberal/q16_fig1.tex}

\begin{solution}
\begin{enumerate}
\item 设正方形 \(ABCD\) 的边长为 \(a\)，设 \(O=AC\cap BD\). 由于 \(ABCD\) 是正方形，得 \(AC\perp BD\).
又因 \(PD\perp\) 底面 \(ABCD\)，所以 \(PD\perp AC\).
在平面 \(PDB\) 内过 \(O\) 作直线 \(l\parallel PD\)，则 \(AC\perp l\) 且 \(AC\perp OB\).
直线 \(l\) 与 \(OB\) 相交于 \(O\)，因此 \(AC\perp\) 平面 \(PDB\).
又因为平面 \(AEC\) 含有直线 \(AC\)，所以平面 \(AEC\perp\) 平面 \(PDB\).

\item 因 \(E\) 是 \(PB\) 的中点，\(O\) 是 \(DB\) 的中点，在三角形 \(PDB\) 中有 \(OE\parallel DP\)，且 \(OE=\frac{1}{2}PD\).
由第一问知 \(AC\perp\) 平面 \(PDB\)，且 \(O\) 是 \(AC\) 与平面 \(PDB\) 的交点，所以 \(AO\perp\) 平面 \(PDB\).
因此 \(OE\) 是 \(AE\) 在平面 \(PDB\) 内的射影，\(\angle AEO\) 就是 \(AE\) 与平面 \(PDB\) 所成的角.

正方形对角线互相平分，所以 \(AO=\frac{1}{2}AC=\frac{\sqrt{2}}{2}AB\).
当 \(PD=\sqrt{2}AB\) 时，\(OE=\frac{1}{2}PD=\frac{\sqrt{2}}{2}AB=AO\).
三角形 \(AOE\) 在 \(O\) 处直角，且 \(AO=OE\)，所以 \(\angle AEO=45^{\circ}\).
故 \(AE\) 与平面 \(PDB\) 所成的角为 \(45^{\circ}\).
\end{enumerate}
\end{solution}
\end{problem}

\begin{problem}
某学生在上学路上要经过 4 个路口，假设在各路口是否遇到红灯是相互独立的，遇到红灯的概率都是 \(\frac{1}{3}\)，遇到红灯时停留的时间都是 \(2 \mathrm{~min}\).
\begin{enumerate}
\item 求这名学生在上学路上到第三个路口时首次遇到红灯的概率；
\item 这名学生在上学路上因遇到红灯停留的总时间至多是 \(4 \mathrm{~min}\) 的概率.
\end{enumerate}

\begin{solution}
\begin{enumerate}
\item 设事件 \(A\) 表示这名学生到第三个路口时首次遇到红灯.
事件 \(A\) 发生当且仅当前两个路口没遇到红灯、第三个路口遇到红灯.
由各路口事件相互独立，
\[
P(A)=\left(1-\frac{1}{3}\right)^2\cdot\frac{1}{3}=\frac{4}{27}
\]

\item 每次遇到红灯停留 \(2\mathrm{~min}\)，设四个路口中遇到红灯的次数为 \(X\)，则停留总时间为 \(2X\mathrm{~min}\).
因此，停留总时间至多为 \(4\mathrm{~min}\) 等价于 \(X\leqslant 2\).
当 \(k=0\)，\(1\)，\(2\) 时，四个路口中恰好有 \(k\) 个遇到红灯的概率为
\(\mathrm C_4^k\left(\frac{1}{3}\right)^k\left(\frac{2}{3}\right)^{4-k}\).
所以
\[
\begin{aligned}
P(X\leqslant 2)
&=\left(\frac{2}{3}\right)^4+\mathrm C_4^1\frac{1}{3}\left(\frac{2}{3}\right)^3+\mathrm C_4^2\left(\frac{1}{3}\right)^2\left(\frac{2}{3}\right)^2\\
&=\frac{16}{81}+\frac{32}{81}+\frac{24}{81}=\frac{8}{9}.
\end{aligned}
\]
\end{enumerate}
\end{solution}
\end{problem}

\begin{problem}
设函数 \( f(x) = x^3 - 3ax + b \)（\(a \neq 0\)）.
\begin{enumerate}
\item 若曲线 \( y = f(x) \) 在点 \((2, f(2))\) 处与直线 \( y = 8 \) 相切，求 \(a\)，\(b\) 的值；
\item 求函数 \( f(x) \) 的单调区间与极值点.
\end{enumerate}

\begin{solution}
\begin{enumerate}
\item \(f'(x)=3x^2-3a\).
直线 \(y=8\) 是水平直线，曲线在点 \((2,f(2))\) 处与它相切，因此必有 \(f'(2)=0\) 且 \(f(2)=8\).
于是
\[
f'(2)=12-3a=0
\]
所以 \(a=4\).
又
\[
f(2)=2^3-3\cdot2a+b=8-6a+b=8
\]
代入 \(a=4\)，得 \(8-24+b=8\)，所以 \(b=24\).

\item
\[
f'(x)=3x^2-3a=3\left(x^2-a\right)
\]
由于 \(a\neq0\)，分两种情形讨论.

当 \(a<0\) 时，\(x^2-a=x^2+(-a)>0\) 对任意 \(x\in\R\) 都成立，所以 \(f'(x)>0\).
此时 \(f(x)\) 在 \(\R\) 上单调递增，没有极值点.

当 \(a>0\) 时，\(f'(x)=0\) 的解为 \(x=\pm\sqrt a\).
当 \(x\in(-\infty,-\sqrt a)\cup(\sqrt a,+\infty)\) 时，\(x^2>a\)，所以 \(f'(x)>0\)，\(f(x)\) 单调递增；当 \(x\in(-\sqrt a,\sqrt a)\) 时，\(x^2<a\)，所以 \(f'(x)<0\)，\(f(x)\) 单调递减.

在 \(x=-\sqrt a\) 处，\(f'(x)\) 由正变负，故为极大值点；在 \(x=\sqrt a\) 处，\(f'(x)\) 由负变正，故为极小值点.
相应的函数值为
\[
f(-\sqrt a)=(-\sqrt a)^3-3a(-\sqrt a)+b=b+2a\sqrt a=b+2a^{\frac{3}{2}}
\]
\[
f(\sqrt a)=(\sqrt a)^3-3a\sqrt a+b=b-2a\sqrt a=b-2a^{\frac{3}{2}}
\]
因此，当 \(a>0\) 时，极大值点为 \(\left(-\sqrt a,b+2a^{\frac{3}{2}}\right)\)，极小值点为 \(\left(\sqrt a,b-2a^{\frac{3}{2}}\right)\).
\end{enumerate}
\end{solution}
\end{problem}

\begin{problem}
已知双曲线 \(C: \frac{x^2}{a^2} - \frac{y^2}{b^2} = 1\)（\(a > 0\)，\(b > 0\)）的离心率为 \(\sqrt{3}\)，右准线方程为 \(x = \frac{\sqrt{3}}{3}\).
\begin{enumerate}
\item 求双曲线 \(C\) 的方程；
\item 已知直线 \(x - y + m = 0\) 与双曲线 \(C\) 交于不同的两点 \(A\)，\(B\)，且线段 \(AB\) 的中点在圆 \(x^{2} + y^{2} = 5\) 上，求 \(m\) 的值.
\end{enumerate}

\begin{solution}
\begin{enumerate}
\item 设双曲线的半焦距为 \(c\)，则 \(c^2=a^2+b^2\)，且离心率 \(e=\frac{c}{a}=\sqrt{3}\).
右准线方程为 \(x=\frac{a^2}{c}=\frac{a}{e}=\frac{\sqrt{3}}{3}a\).
由题设 \(\frac{\sqrt{3}}{3}a=\frac{\sqrt{3}}{3}\)，得 \(a=1\).
于是 \(c=\sqrt{3}\)，\(b^2=c^2-a^2=3-1=2\)，又因 \(b>0\)，得 \(b=\sqrt{2}\).
所以双曲线 \(C\) 的方程为
\[
x^2-\frac{y^2}{2}=1
\]

\item 由直线方程 \(x-y+m=0\)，得 \(y=x+m\).
代入双曲线方程，得
\[
x^2-\frac{(x+m)^2}{2}=1
\]
整理得
\[
x^2-2mx-m^2-2=0
\]
其判别式为
\[
\Delta=(-2m)^2-4(-m^2-2)=8(m^2+1)>0
\]
所以对任意实数 \(m\)，该直线都与双曲线交于两个不同点.

设这两个交点的坐标分别为 \(A(x_1,y_1)\)，\(B(x_2,y_2)\).
由一元二次方程的根与系数关系，\(x_1+x_2=2m\).
又 \(y_1=x_1+m\)，\(y_2=x_2+m\)，所以线段 \(AB\) 的中点为
\[
M\left(\frac{x_1+x_2}{2},\frac{y_1+y_2}{2}\right)=\left(m,2m\right)
\]
由中点 \(M\) 在圆 \(x^2+y^2=5\) 上，得
\(m^2+(2m)^2=5\)，\(\Longrightarrow\)，\(5m^2=5\)，\(\Longrightarrow\)，\(m=\pm1\)
故 \(m\) 的值为 \(1\) 或 \(-1\).
\end{enumerate}
\end{solution}
\end{problem}

\begin{problem}
设数列 \(\{a_{n}\}\) 的通项公式为 \(a_{n} = pn + q\)（\(n \in \N^{*}\)，\(p > 0\)）. 数列 \(\{b_{m}\}\) 定义如下：对于正整数 \(m\)，\(b_{m}\) 是使得不等式 \(a_{n} \geqslant m\) 成立的所有 \(n\) 中的最小值.
\begin{enumerate}
\item 若 \(p = \frac{1}{2}\)，\(q = -\frac{1}{3}\)，求 \(b_{3}\)；
\item 若 \(p = 2\)，\(q = -1\)，求数列 \(\{b_m\}\) 的前 \(2m\) 项和公式；
\item 是否存在 \(p\) 和 \(q\)，使得 \(b_{m} = 3m + 2\)（\(m \in \N^{*}\)）？ 如果存在，求 \(p\) 和 \(q\) 的取值范围；如果不存在，请说明理由.
\end{enumerate}

\begin{solution}
\begin{enumerate}
\item 当 \(p=\frac{1}{2}\)，\(q=-\frac{1}{3}\) 时，
\(a_n=\frac{1}{2}n-\frac{1}{3}\).
由 \(a_n\geqslant3\)，得
\[
\frac{1}{2}n-\frac{1}{3}\geqslant3
\quad\Longleftrightarrow\quad
n\geqslant\frac{20}{3}
\]
满足条件的最小正整数为 \(7\)，所以 \(b_3=7\).

\item 当 \(p=2\)，\(q=-1\) 时，\(a_n=2n-1\).
对于正整数 \(m\)，
\[
2n-1\geqslant m
\quad\Longleftrightarrow\quad
n\geqslant\frac{m+1}{2}
\]
对任意 \(k\in\N^*\)，当 \(m=2k-1\) 时，\(b_{2k-1}=k\)；当 \(m=2k\) 时，\(b_{2k}=k+1\).
因此前 \(2m\) 项的和为
\[
\begin{aligned}
\sum_{j=1}^{2m}b_j
&=\sum_{k=1}^{m}\left(b_{2k-1}+b_{2k}\right)
=\sum_{k=1}^{m}(2k+1)\\
&=2\cdot\frac{m(m+1)}{2}+m=m^2+2m.
\end{aligned}
\]
所以数列 \(\{b_m\}\) 的前 \(2m\) 项和公式为 \(m^2+2m\).

\item 因为 \(p>0\)，\(a_n=pn+q\) 随 \(n\) 严格递增.
若要使得 \(b_m=3m+2\)，则对每个 \(m\in\N^*\)，必须同时满足
\[
p(3m+2)+q\geqslant m
\]
和
\[
p(3m+1)+q<m
\]
前一个不等式表示 \(n=3m+2\) 已经使不等式成立，后一个不等式表示前一个正整数 \(n=3m+1\) 还不能使不等式成立.
将上述不等式整理为
\((3p-1)m+2p+q\geqslant0\)，\((3p-1)m+p+q<0\).
这两个不等式对所有正整数 \(m\) 都成立.
如果 \(3p-1>0\)，则第二个不等式当 \(m\) 足够大时不成立；如果 \(3p-1<0\)，则第一个不等式当 \(m\) 足够大时不成立.
所以必有 \(3p-1=0\)，即 \(p=\frac{1}{3}\).
此时上述两个不等式分别变为
\(\frac{2}{3}+q\geqslant0\)，\(\frac{1}{3}+q<0\).
因此
\[
-\frac{2}{3}\leqslant q< -\frac{1}{3}
\]
反过来，当 \(p=\frac{1}{3}\) 且 \(-\frac{2}{3}\leqslant q< -\frac{1}{3}\) 时，有
\(a_{3m+2}=m+\frac{2}{3}+q\geqslant m\)，\(a_{3m+1}=m+\frac{1}{3}+q<m\).
由 \(a_n\) 严格递增，可知 \(b_m=3m+2\).
所以存在这样的 \(p\) 和 \(q\)，其取值范围为
\(p=\frac{1}{3}\)，\(-\frac{2}{3}\leqslant q< -\frac{1}{3}\).
\end{enumerate}
\end{solution}
\end{problem}
